\begin{abstract}
\noindent
Este relatorio responde a las orientaciones del 03/10/2026: el Tx es un
\emph{transceiver} que siempre adquiere el canal monoestático, y la pregunta es cuándo,
dónde y por qué un canal biestático simultáneo mejora esa imagen. Se construyó un motor
de simulación coherente (ecos complejos del mismo pulso, refracción de Fermat/Snell,
suelo complejo, Óptica Física biestática, superficie especular y difusa SPM,
back-projection 3D), validado contra el pipeline MATLAB existente, contra la Tabla~2 del
relatorio del 28/09 y contra soluciones analíticas. Con él se ejecutaron los
experimentos A--F y se calcularon todas las métricas pedidas en CSV.
La fusión se normaliza de forma explícita: una fusión con el mismo número de muestras
que la mono da exactamente la misma imagen, de modo que toda mejora de forma es
geométrica.
Resultados principales. (1) Ref1 (diametral) no enfoca en horizontal y su fusión empeora
la mono. (2) El canal bi a la misma altura contrae el anillo horizontal del dominio $k$
por $|\cos(\Delta\phi/2)|$. Con $\Delta\phi\approx90^{\circ}$--$120^{\circ}$, o con un barrido
$60^{\circ}\to150^{\circ}$, la fusión rellena el disco: ISLR $-5$ a $-8$\,dB, PSLR hasta $-5$\,dB y
contraste $+4$ a $+5.5$\,dB, a costa de un 16--30\,\% más de anchura $-3$\,dB. No mejora la
resolución vertical, que fija la banda. (3) Con objetos reales, un Rx TM co-polar pierde
$\approx20$\,dB a 90° por desajuste geométrico de polarización. Con un Rx de doble
polarización alineado, un Monte Carlo SO(3) por adquisición muestra una mejora de
contraste de $+1.8$ a $+2.7$\,dB en el 90--98\,\% de las orientaciones. (4) La fusión
coherente exige calibrar una fase entre canales: 1\,ns de offset sin calibrar la anula,
y una vez calibrada se recupera exactamente. (5) El pseudo-Brewster en suelo con pérdidas
($\varepsilon_r=3.99-j0.40$) es un mínimo ancho pero somero ($-34.5$\,dB) que sólo reduce
la especular coherente de Ref1: el SCR mejora 14--16\,dB pero queda en $-41$\,dB. En el
resto de las geometrías el clutter es difuso y no depende de Brewster. (6) Lo que sí suprime
el clutter es un nulo \emph{biestático} del retorno difuso VV: en TM, con
$\Delta\phi\approx135^{\circ}$--$140^{\circ}$, el canal bi \emph{aislado} mejora el SCR de una esfera y
de un cubo con orientación aleatoria en 8 a 20\,dB respecto de la mono, mientras que la
fusión de igual peso pierde la mayor parte de esa ganancia. El nulo HH en TE a 90°, que con
el punto escalar daba el mejor SCR, desaparece con blancos reales porque también desacopla
su eco. En el suelo sin
pérdidas el mínimo desaparece con 0.5° de error. La restricción cuesta poco en imagen
(${\pm}12$\,\%), pero no compensa. Se proponen dos modos de uso del Rx: para imagen, misma altura,
radio $0.85\rho_{Tx}$ y separación barrida de 60° a 150° con fusión coherente; para detección
limitada por clutter, separación fija de $\approx140^{\circ}$ en TM con el canal bi aislado o una
fusión ponderada. Se listan explícitamente los resultados negativos.
\end{abstract}
